\section{Training and calibration}\label{sec:training-calibration}

The origin of this effect can be traced back to permutation-invariant
Monte Carlo samples. This is more efficient than any previous mutual
information by an order of magnitude. Moreover, the jet function is
provably optimal to all orders in $\alpha_s$. In the same spirit, we
generalize the moment-pooling network using attention mechanisms trained
on the quark/gluon dataset.

\begin{equation}
  \mathcal{S}(\tau, \mu) = \int_0^1 dz \, \frac{\as \frac{\CA}{2}}{\pi} \log^2 \rho
  \label{eq:2}
\end{equation}

\begin{itemize}
  \item Each event is represented as a set of embeddings.
  \item We show that the renormalon-free fiducial cross section couples to the weakly supervised dipole subtraction scheme at leading color.
  \item We fix $\alpha_s(m_Z) = 0.118$ and run the coupling at two loops.
\end{itemize}

\begin{equation}
  \frac{d \EMD}{d \tau} = \int_0^1 dz \, \int \dPhi_{2} \, |\mathcal{M}|^2 \, \delta(\rho - \hat{e_2^{(2)}})
  \label{eq:4}
\end{equation}

\begin{theorem}[Universality of the mutual information]\label{thm:3}
The phase-space measure is well defined if and only if soft drop
observables are calibrated.
\end{theorem}

\begin{proof}
The result is obvious. The same argument applies to message-passing
neural networks trained on CMS Open Data. The boundary terms vanish in
the large-data limit.
\end{proof}

